\documentclass[12pt, a4paper]{article}

\usepackage{cmap}
\usepackage[utf8]{inputenc}
\usepackage[T2A]{fontenc}
\usepackage[english,russian]{babel}
\usepackage{authblk}

\usepackage{mathtools}
\usepackage{amssymb}
\usepackage{amsthm}
\usepackage{amsmath}


\begin{document}
	
	\section*{Формула трапеций (запись через веса)}
	Вещественная сетка из $N$ узлов ($N - 1$ шагов).
	\begin{equation}
		t_{n} = t_{\text{start}} + h_{t}n, n \in \left[0, N - 1\right], h_{t} = \left(t_{\text{finish}} - t_{\text{start}}\right) / \left(N - 1\right).
	\end{equation}
	Комплексная сетка из $M$ узлов ($M - 1$ шагов).
	\begin{equation}
		z\left(s\right) = \varphi\left(s\right), h_{z}\left(s\right)=z\left(s+h_{s}\right) - z\left(s\right) \approx \varphi'\left(s\right)h_{s}.
	\end{equation}
	\begin{equation}
		z_{m} = z\left(s_{m}\right) = \varphi\left(s_{m}\right), h_{z}\left(s_{m}\right)=z\left(s_{m+1}\right) - z\left(s_m\right) \approx \varphi'\left(s_{m}\right)h_{s}.
	\end{equation}
%	\begin{equation}
%		\varphi\left(s\right) = z_{0} + R\exp\left(2\pi i s\right), \varphi'\left(s\right) = 2\pi i R\exp\left(2\pi i s\right).
%	\end{equation}
	\begin{multline}
		s_{m} = s_{\text{start}} + h_{s}m, m \in \left[0, M - 1\right],\\
		 h_{s} = \left(s_{\text{finish}} - s_{\text{start}}\right) / \left(M - 1\right), s_{\text{start}} = 0, s_{\text{finish}} = 1.
	\end{multline}
%	\begin{multline}
%		z_{m} = z_{0} + R\exp\left(2\pi i s_{m}\right), m \in \left[0, M - 1\right],\\
%		 h_{z} = 2\pi i h_{s} R\exp\left(2\pi i s_{m}\right)=
%		 2\pi i h_{s}\left(z_{m}-z_{0}\right), m\in\left[0, M-1\right].
%	\end{multline}
%	Интеграл Коши
%	\begin{equation}
%		u\left(t\right) = \frac{1}{2\pi i}\int_{\left|z\right|=1}\frac{\tilde{u}\left(z\right)}{z-t}dz.
%	\end{equation}
	Интеграл по вещественной оси
	\begin{equation}
		I=\int_{t_s}^{t_f}u\left(t\right)dt.
	\end{equation}
	Формула трапеций на вещественной сетке (общий случай)
	\begin{multline}
		I_{N} = \sum_{n = 0}^{N - 2}\frac{u_{n+1} + u_{n}}{2}\left(t_{n+1} - t_{n}\right) =\\
		=\sum_{n = 0}^{N - 1}g_{n}u_{n},\quad
		g_{n}=
		\left\{\begin{aligned}
			&\left(t_{1} - t_{0}\right)/2, &&n = 0,\\
			&\left(t_{n+1} - t_{n-1}\right)/2, &&n\in\left[1,N-2\right],\\
			&\left(t_{N-1} - t_{N-2}\right)/2, &&n=N-1.
		\end{aligned}\right.
	\end{multline}
	На равномерной вещественной сетке
	\begin{equation}
		I_{N} = h_{t}\sum_{n = 0}^{N - 1}g_{n}u_{n},\quad
		g_{n}=
		\left\{\begin{aligned}
			&1/2, &&n = 0,\\
			&1, &&n\in\left[1,N-2\right],\\
			&1/2, &&n=N-1.
		\end{aligned}\right.
	\end{equation}
	Интеграл по контуру $C$ в комплексной плоскости
	\begin{equation}
		\tilde{I}=\int_{C}\tilde{u}\left(z\right)dz=\int_{0}^{1}\tilde{u}\left(\varphi\left(s\right)\right)\varphi'\left(s\right)ds.
	\end{equation}
	Предположим, что контур гладкий и параметризуется функцией $\varphi\left(s\right),~s\in\left[0,1\right]$. Тогда формула трапеций примет следующий вид (общий случай)
	\begin{multline}
		\tilde{I} = h_{s}\sum_{m = 0}^{M - 2}\frac{\tilde{u}_{m+1}\varphi'_{m+1} + \tilde{u}_{m}\varphi'_{m}}{2} =h_{s}\sum_{m=0}^{M-1}\tilde{g}_{m}\tilde{u}_{m},\\
		\tilde{g}_{m}=
		\left\{\begin{aligned}
			&\varphi'_{0}/2, &&m = 0,\\
			&\varphi'_{m}, && m\in\left[1,M-2\right],\\
			&\varphi'_{M-1}/2 && m = M-1.
		\end{aligned}\right.
	\end{multline}
	То же, но на замкнутом контуре ($\tilde{u}_{M-1} = \tilde{u}_{0}$)
	\begin{equation}
		\tilde{I} =h_{s}\sum_{m=0}^{M-2}\tilde{g}_{m}\tilde{u}_{m},~
		\tilde{g}_{m}=
		\left\{\begin{aligned}
			&\frac{\varphi'_{0} + \varphi'_{M-1}}{2}=\varphi'_{0}, &&m = 0,\\
			&\varphi'_{m}, && m\in\left[1,M-2\right].
		\end{aligned}\right.
	\end{equation}
	
%	\section*{Вычисление контурного интеграла}
%	Исходный интеграл
%	
%	Квадратура трапеций
%	\begin{equation}
%		\tilde{I}_{M}= h_{s}\sum_{m=0}^{M-2}\frac{\tilde{u}\left(\varphi\left(s_{m+1}\right)\right)\varphi'\left(s_{m+1}\right)+
%		\tilde{u}\left(\varphi\left(s_{m}\right)\right)\varphi'\left(s_{m}\right)}{2}.
%	\end{equation}
	
	
	
\end{document}